\section{Statistical design}
\label{sec:stats}

\paragraph{Unit of analysis and primary outcome.}
The unit of analysis is the independent stochastic replicate; seeds are derived as
\texttt{SeedSequence([base\_seed, point, replicate])}, and common random numbers are used across design points in the
screening and variance-based designs. Each run is scored by the automatic evaluator of
Section~\ref{sec:observations}. The primary per-run outcome is \emph{O17p}: the run is Calhoun-like, meaning that
every applicable essential pattern passes and no anti-pattern is triggered, \emph{and} its isolated class of O13 is
persistent. With the default initial condition ($N_0=400$), O1--O3 are not applicable and are excluded from O17p;
they are studied separately with the founder protocol. Every other essential pattern must pass: a pattern that
cannot be evaluated in a run (NA) counts as a failure of O17p, so that no pattern is passed vacuously.
The primary evaluator, which is not preregistered, differs from the preregistered one: it takes the peak time at the start of the plateau ($t'_{\rm pk}$, first
week within 2\% of the maximum), the end of natality at 99\% of the births ($t_{B99}$), O4 with free space at the
peak, O10 on the adults, O11 on the cohorts born in the second half of the natality and the crowded class at $m>8$;
every run is also scored with the preregistered evaluator and with each option changed separately
(Appendix~\ref{app:res-rescoring}). Some patterns are defined for an ensemble and scored per run: the per-run
version of O7 is ``extinct before the horizon $T$'', its ensemble version is $P_{\rm ext}(T)\ge0.9$ estimated by
Kaplan--Meier, and the fraction of runs that pass O7 is the binomial estimate of the latter.

\paragraph{Presets and stages.}
All designs are expressed relative to the reference candidate \textsf{C1} (Table~\ref{tab:params}). The controls
are \textsf{C0} (\textsf{C1} without R7), \textsf{C1} with $\alpha=0$, and the short-lived and long-lived controls
(the core model with the four rules switched off, $\alpha=1.4$, and $a_{\max}=60$ or 120;
Section~\ref{sec:model-candidates}); the confirmatory contrast is \textsf{C2}; the null models
(Section~\ref{sec:res-null}) are run from the initial condition of \textsf{C1}. Rule ablations switch a rule off by
setting its parameters to their neutral values (for R2, full inheritance and individual recovery at the core
rate $r=0.01$). The design has six kinds of stages: production ensembles of the candidates, controls and null models, with
one reference ensemble for \textsf{C1}, \textsf{C1}$'$, \textsf{C0} and the $\alpha=0$ control; a $2^4$ factorial
ablation with implementation variants; transplant experiments (paired cross-transplant, mixed transplants with
healthy partners, sampling at a common absolute time, transplants of adults of intermediate social state); phase
planes in the socialisation, space and refuge parameters, a seeding-density plane and a founder plane; robustness
variants (background mortality, nest period, asynchronous update, pups born above the deterioration threshold,
non-synchronous initial condition, lattice size at fixed density); and global sensitivity analysis in three steps:
Morris elementary effects over 15 factors (20 trajectories, 10 runs per point), two independent 15-factor Latin
hypercubes (200 points, 10 runs each, plus 40 runs at the points near the boundary), and Sobol' indices with the
Saltelli design and Jansen estimators \cite{Saltelli2010} for 11 factors (the 8 retained by the screening plus $w$,
$r$ and $a_w$), with $N=256$ and 3 runs per point, extended to 6 \cite{tenBroeke2016}.
Morris and Sobol' indices are descriptive, not tests; they are reported with bootstrap intervals and with an
\emph{empirical noise floor}, the same estimators applied to the split-half contrast of the replicate means at each
design point (Appendix~\ref{sec:res-sensitivity}).

\paragraph{Confirmatory hypotheses.}
Seven hypotheses were fixed before the production runs, each with a quantitative prediction derived from the
mean-field analysis or the calibration. The plan, with designs and predictions, was fixed on 1 October 2026, 22:02 UTC, before the launch of the production
runs, and is reproduced in \texttt{docs/preregistered\_plan.md} of the code repository; the primary evaluator, the stages not in that plan and
their decision rules are not preregistered, but were fixed (\texttt{docs/production\_plan.md}, 2 October 2026, 05:34 UTC) after exploratory runs
with the primary evaluator (smoke tests on the data of the preregistered design and on two replicates per point) and
before the production runs. One option of the primary evaluator, the birth reference $t_{B99}$ of O5, was chosen after
those exploratory runs had shown that the alternative fails O5 in most \textsf{C1} runs; it is the only option on
which a confirmatory decision (Q2) depends, and both are reported (Table~\ref{tab:deviations}).
They form one family with Holm correction at $\alpha=0.05$: Q1, structural irreversibility (the same phase-D
animals transplanted to a \textsf{C1} enclosure and to one without R1 and R2; exact one-sided McNemar test); Q2, each
rule is necessary (intersection--union test over the four single-rule ablations, one-sided Fisher tests, $n=120$ per
arm); Q3, near-deterministic collapse versus the short-lived control (Fisher test on $P_{\rm ext}$ and the coefficient of
variation of $T_{\rm ext}$); Q4, the phase-D law, $T_{\rm ext}-t_{\rm last}$ one-to-one with $a_{\max}$ (TOST with
margins $[0.8,1.25]$); Q5, the coexistence band of O13 centred at $n^*\approx m_c$ (Firth regression on $\log n^*$);
Q6, the plasticity window weighs less than the learning rate (profile likelihood of $\theta$ in
$\log r+\theta\log a_w$, $\theta<1$; quasi-binomial $p$); Q7, endogenous collapse without crowding damage in the
subcritical corner (Fisher).
A second confirmatory family compares \textsf{C1} and \textsf{C2} on every pattern and on six continuous magnitudes
($n=200$ per arm; two-sided Fisher and Mann--Whitney tests, log-rank for $T_{\rm ext}$, Holm correction over 29 rows,
which hold 27 distinct tests because O6, O7 and A2 are decided by the same four \textsf{C2} colonies; nine rows are constant in both arms
and enter with $p=1$, so the correction is conservative). One confirmatory test is not preregistered but was fixed before the production runs: the equivalence of the exponent of $a_w$ with that of $r$ at $\gamma=0$ (TOST with margins
$[2/3,3/2]$), with a likelihood ratio for $\theta=1$; its design turned out to lack the resolution for that test
(the logistic model is separated), so the stage is reported non-parametrically (Section~\ref{sec:res-Q}). Rejection
of a null hypothesis and fulfilment of the quantitative prediction are
reported separately. A Holm-adjusted $p$ rejects if it is at most the level, and adjusted values within 0.01 of
the level are given with four digits. Q2 is an intersection--union test: it rejects only if every component test
rejects at the full level, so the four component tests are not corrected among themselves. Exploratory contrasts
are grouped into families and controlled with Benjamini--Hochberg at 10\%. The rules for the stages not in the preregistered plan were also fixed before the runs: a criterion is \emph{generic} if it passes in at least two of the four null models, a
robustness variant \emph{preserves} the result if O17p $\ge0.8$, and the feasible fraction of the hypercubes is
reported with the empirical-Bayes estimator of Appendix~\ref{sec:res-sensitivity}.

\paragraph{Intervals, survival and budget.}
Unless stated otherwise, intervals are 95\% intervals: Wilson for proportions; Newcombe for differences of two
proportions; percentile bootstrap, resampling runs or design points, for the coefficient of variation and the
spread of extinction times, for the Calhoun-like fraction of the hypercubes and for sensitivity indices; profile
likelihood for the exponents $\hat\theta$, quasi-binomial when stated; Greenwood log--log for Kaplan--Meier curves
and medians. Extinction times are right-censored at the horizon $T$, and means over extinct runs only are never
reported: we use Kaplan--Meier estimates, $P_{\rm ext}(T)$ with Wilson intervals, log-rank comparisons, Weibull fits
and a cure model with an escaping fraction where a fraction of the colonies survives \cite{Ovaskainen2010}; the
spread of the extinction time is the interquartile range over the median of the Kaplan--Meier estimate.
The 95\% Wilson half-width at $p=0.5$ is $\pm0.17$, $\pm0.12$, $\pm0.09$ and $\pm0.07$ for $n=30$, $60$, $120$ and
$200$; production, contrast, null-model and robustness ensembles have 200 runs (100 for the short-lived and long-lived
controls, 60 for the background-mortality curve and the lattice sizes), single ablations 120, multiple ablations 40,
transplant arms 60 replicates, phase-plane points 30 (100 at the 40 points nearest to a boundary), Morris and
hypercube points 10 \cite{Lorscheid2012}. No power analysis was done after the runs: observed
power is a monotone function of the $p$-value and cannot explain a non-rejection \cite{HoenigHeisey2001}, so we
report effect sizes with their intervals instead. The production program comprises 51\,828 simulations and 960
transplant replicates in 26 stages, about 4.4 hours of logged wall time on four cores, 1.8 of them for the Sobol'
design (Table~\ref{tab:res-ensembles}); a run of \textsf{C1} or one of its variants costs 0.5--2.6\,s of CPU time.
Table~\ref{tab:deviations} (Appendix~\ref{app:res-deviations}) lists every deviation from the
preregistered plan, including the redefinition of the primary evaluator; analyses not fixed before the production runs are labelled
post hoc wherever they are reported.
